% ---------------------------------------------------------------------------
% Self-contained file: paste this WHOLE file into Overleaf's main.tex
% (replace everything that is currently there).
%
% IMPORTANT: in Overleaf, set the compiler to XeLaTeX before recompiling:
%   Menu (top-left) -> Compiler -> XeLaTeX -> Recompile
% pdfLaTeX (the Overleaf default) cannot render the Vietnamese diacritics
% used in one example sentence below (u+ hook, o+ hook, combining tone
% marks); XeLaTeX + fontspec + a Unicode font (Noto Serif, bundled with
% every TeX Live install Overleaf uses) renders them correctly.
%
% If you already have your own main.tex with its own chapters, don't paste
% this whole file -- just copy the \chapter{...}...\end{document}-less body
% from report_src/chapter4_implementation.tex into your own document after
% Chapter 3, and make sure your own preamble also loads amsmath, booktabs,
% hyperref, and (if you keep the Vietnamese example) fontspec with a
% Unicode font under XeLaTeX/LuaLaTeX.
% ---------------------------------------------------------------------------
\documentclass[11pt]{report}

\usepackage{fontspec}
\setmainfont{Noto Serif} % full Unicode/Vietnamese coverage; change if you use a different font elsewhere in the report

\usepackage{amsmath}     % equation/eqref
\usepackage{booktabs}    % \toprule \midrule \bottomrule
\usepackage{hyperref}    % \url{} and clickable cross-references
\usepackage[margin=1in]{geometry}

% So the first \chapter below is numbered "Chapter 4" (continuing after
% Chapters 1-3 of the existing report). Remove this line if you paste the
% body into a document that already has Chapters 1-3 typeset before it.
\setcounter{chapter}{3}

\begin{document}

\chapter{Experiments and Results}
\label{chap:experiments}

\section{Implementation}
\label{sec:implementation}

This section describes how the Graph-RAG pipeline of Chapter~3 is realised
as a working system: the concrete data artefacts consumed at run time
(Section~\ref{sec:data}), the software architecture
(Section~\ref{sec:architecture}), the retrieval and ranking parameters
(Section~\ref{sec:retrieval-impl}), how the large language model (LLM) is
invoked and prompted (Section~\ref{sec:llm-usage}), and how its output is
validated before being shown to a user (Section~\ref{sec:grounding}).

\subsection{Data}
\label{sec:data}

The knowledge graph is built directly from the Decagon dataset, distributed
by the Stanford SNAP BioSNAP collection\footnote{\url{https://snap.stanford.edu/decagon/}}
as four plain CSV files; the fourth, \texttt{bio-decagon-mono.csv}, is
optional and is used only as a ranking feature (Section~\ref{sec:retrieval-impl}),
never for retrieval. Table~\ref{tab:data-summary} reports the exact size of
each file as distributed -- obtained by loading all four files with no
filtering -- and matches the totals already summarised in Section~1.3.2 up
to the small differences expected from counting raw rows rather than
deduplicated entities.

\begin{table}[t]
  \caption{Decagon dataset, as distributed (unfiltered row and entity counts)}
  \label{tab:data-summary}
  \centering
  \begin{tabular}{@{}lrl@{}}
    \toprule
    File & Rows & Key entities \\
    \midrule
    \texttt{bio-decagon-targets.csv}         & 18{,}690    & 284 drugs, 3{,}648 proteins \\
    \texttt{bio-decagon-ppi.csv}             & 715{,}612   & 19{,}081 proteins \\
    \texttt{bio-decagon-combo.csv}           & 4{,}649{,}441 & 645 drugs, 63{,}473 drug pairs, 1{,}317 effect codes \\
    \texttt{bio-decagon-mono.csv} (optional) & 174{,}977   & 639 drugs, 10{,}184 effect codes \\
    \bottomrule
  \end{tabular}
\end{table}

Two structural properties of the graph directly motivated implementation
choices in Section~\ref{sec:retrieval-impl}. First, both the drug--target
and the protein--protein subgraphs are heavily right-skewed
(Table~\ref{tab:degree}): half of all drugs bind at most 11 proteins, yet
one outlier binds 2{,}840, and half of all proteins have at most 33
interaction partners, yet the busiest protein has 2{,}519. Treating every
protein identically during bridge search would let a handful of such hub
proteins connect almost any two drugs through a biologically uninformative
path. Fixing \texttt{hub\_degree\_threshold} at 500 (Table~\ref{tab:bounds})
removes only the extreme tail -- 220 of 19{,}081 proteins (1.15\%), close to
the 99th percentile of the observed degree distribution -- while leaving the
rest of the PPI network intact.

\begin{table}[t]
  \caption{Degree-distribution statistics of the drug--target and protein--protein graphs}
  \label{tab:degree}
  \centering
  \begin{tabular}{@{}lrrrrrr@{}}
    \toprule
    Degree & Min & Median & Mean & P90 & P99 & Max \\
    \midrule
    Targets per drug            & 1 & 11 & 65.8 & 170 & 206 & 2{,}840 \\
    Drugs per target protein    & 1 & 1  & 5.1  & 5   & 78  & 94 \\
    PPI neighbours per protein  & 1 & 33 & 75.0 & 201 & 525 & 2{,}519 \\
    \bottomrule
  \end{tabular}
\end{table}

Second, the 645 drugs that carry documented pairwise history in
\texttt{bio-decagon-combo.csv} are \emph{not} the same population as the 284
drugs listed in \texttt{bio-decagon-targets.csv}: only 284 of the 645 (44\%)
have a known target set. For the remaining 361 drugs, Tier-1 evidence (a
documented pairwise effect, Table~3.2) is still available whenever the pair
itself was tested, but Tier-2 amplification and Tier-3 bridging are
reachable only for drugs with target annotations, since both require a
protein-level path. This asymmetry is a property of the released data, not
of the retrieval algorithm, and bounds how much mechanistic -- as opposed to
purely historical -- evidence the pipeline can offer for target-unannotated
drugs.

Third, the volume of documented pairwise history is itself heavy-tailed at
two levels (Table~\ref{tab:effect-freq}). A drug pair that has \emph{any}
documented effect has a median of 53 and up to 540 of them, which is why
effects are ranked (Eq.~\eqref{eq:score}) rather than listed in full; and a
given effect code is documented for a median of 1{,}596 pairs but as few as
1, which is why the inverse-frequency term of Eq.~\eqref{eq:idf} matters in
practice. Of the 1{,}317 distinct effect codes present in the raw file, 963
occur in at least 500 pairs -- matching, within one effect, the 964 side
effects that Zitnik et al.~[1] report evaluating after applying the same
frequency cut-off -- while the remaining $\sim$350 codes are comparatively
rare and receive a correspondingly larger $I(e)$ rather than being dropped.

\begin{table}[t]
  \caption{Side-effect frequency statistics in \texttt{bio-decagon-combo.csv}}
  \label{tab:effect-freq}
  \centering
  \begin{tabular}{@{}lrrrrrr@{}}
    \toprule
     & Min & Median & Mean & P90 & P99 & Max \\
    \midrule
    Effects per documented pair      & 1 & 53      & 73.3    & 165   & 305    & 540 \\
    Pairs documenting a given effect & 1 & 1{,}596 & 3{,}530.3 & 9{,}530 & 21{,}396 & 28{,}568 \\
    \bottomrule
  \end{tabular}
\end{table}

These characteristics also shape how the raw files are turned into a
queryable graph at run time. Drug--target and protein--protein interaction
edges are loaded once into in-memory adjacency maps
(\texttt{dict[str, set[str]]}), giving $O(1)$ neighbour look-ups during
retrieval. The combination file, at 4{,}649{,}441 rows, is not re-parsed on
every run: it is materialised once into a local SQLite cache
(\texttt{side\_effects(drug\_a, drug\_b, effect\_code, effect\_name)},
stored \texttt{WITHOUT ROWID} with a composite primary key, plus an index on
\texttt{drug\_b} for symmetric pair look-ups), together with two derived
tables used directly by the ranking formula of
Section~\ref{sec:retrieval-impl}: \texttt{effect\_counts} ($n_e$, the number
of pairs documenting each effect) and \texttt{stats} ($N$, the total number
of distinct documented pairs). A SHA-256 signature of each source file's
size and modification time, together with a cache-schema version, is stored
in a \texttt{metadata} table; the cache is rebuilt automatically only when a
source file or the schema changes, which keeps repeated start-ups fast and
avoids re-reading millions of rows on every process launch.

A free-text drug name typed by a user -- including Vietnamese brand names,
typos, or several drugs written in one sentence -- is resolved to a Decagon
STITCH identifier by a dedicated resolver. Resolution is attempted, in
order, against (i) a STITCH-ID pattern, (ii) an exact match on normalised
generic names, synonyms, or Vietnamese brand aliases, and (iii) fuzzy string
matching (RapidFuzz\footnote{\url{https://github.com/rapidfuzz/RapidFuzz}},
indel ratio), accepted automatically above a score of 90 and otherwise
offered only as a suggestion above 75. A curated stop-word list removes
conversational filler (dosage, timing, and symptom words, in both
Vietnamese and English) before matching, so that a message such as
\emph{"em đang uống panadol với ibuprofen"} resolves to the two intended
drugs rather than being treated as one unmatched phrase. Name resolution
and prompt serialisation draw on five supporting resource tables, built
offline by the \texttt{scripts/} utilities: canonical generic names per
STITCH ID, Vietnamese brand aliases, gene symbols, Vietnamese translations
of effect names, and two curated lists marking EMA Important Medical
Events and life-threatening effects.

\subsection{Software Architecture}
\label{sec:architecture}

The system is implemented in Python. Table~\ref{tab:modules} summarises the
responsibility of each module under \texttt{graphrag/}; \texttt{app.py} is a
Streamlit\footnote{\url{https://streamlit.io}} front-end that calls the
single orchestration entry point, \texttt{PolypharmacySystem.predict()}.
Configuration -- file paths, the LLM endpoint and model, and the retrieval
and ranking parameters of Table~\ref{tab:bounds} -- is read from
environment variables through a \texttt{.env} file, so the retrieval bounds
and the model used for generation can be changed without editing code.
External dependencies are kept deliberately small: \texttt{httpx} for the
synchronous HTTP call to the LLM provider, \texttt{pydantic} for schema
validation of the LLM's output, \texttt{rapidfuzz} for name resolution, and
\texttt{pyvis} for rendering the interaction graph in the user interface.

\begin{table}[t]
  \caption{Module responsibilities}
  \label{tab:modules}
  \centering
  \begin{tabular}{@{}ll@{}}
    \toprule
    Module & Responsibility \\
    \midrule
    \texttt{config.py}            & Environment-driven settings, ranking weights \\
    \texttt{dataset.py}           & Graph loading, SQLite cache, subgraph retrieval \\
    \texttt{names.py} / \texttt{resolver.py} & Name normalisation and drug-name resolution \\
    \texttt{ranking.py}           & Regimen-level effect scoring (Sec.~\ref{sec:retrieval-impl}) \\
    \texttt{llm.py}               & Prompt construction, API call, schema validation \\
    \texttt{presenter.py}         & Assembling evidence into UI-ready structures \\
    \texttt{pipeline.py}          & \texttt{PolypharmacySystem}: end-to-end orchestration \\
    \texttt{visualize.py}         & Interaction-graph rendering (PyVis) \\
    \bottomrule
  \end{tabular}
\end{table}

\subsection{Subgraph Retrieval, Pruning, and Ranking}
\label{sec:retrieval-impl}

Retrieval follows the six steps of Table~3.1, subject to the bounds in
Table~\ref{tab:bounds}. A protein is pruned as a generic hub, and excluded
from any bridge, once its number of PPI neighbours exceeds
\texttt{hub\_degree\_threshold}; a direct target of a queried drug is always
retained regardless of degree. Tier-3 bridge search is bounded per drug
pair (\texttt{tier3\_one\_hop\_per\_pair}, \texttt{tier3\_two\_hop\_per\_pair})
and stops as soon as both caps are full and one further candidate has
confirmed that more bridges exist, so the search is deterministic and never
silently under-reports truncation. The total evidence serialised to the LLM
is capped at \texttt{max\_evidence} and allocated across drug pairs by
round-robin selection, so that a single high-degree pair cannot consume the
whole budget. Documented evidence is built only for the effects that will be
shown, and Tier~2 and Tier~3 each have a reserved quarter of the budget, so
Tier~1 cannot starve them; unused slots go to the other tiers.

\begin{table}[t]
  \caption{Retrieval and pruning bounds (default values)}
  \label{tab:bounds}
  \centering
  \begin{tabular}{@{}lrl@{}}
    \toprule
    Parameter & Default & Role \\
    \midrule
    \texttt{hub\_degree\_threshold}      & 500 & PPI degree above which a protein is pruned \\
    \texttt{max\_evidence}               & 400 & Cap on evidence items serialised to the LLM \\
    \texttt{tier3\_one\_hop\_per\_pair}  & 8   & Max.\ 1-hop bridges sampled per undocumented pair \\
    \texttt{tier3\_two\_hop\_per\_pair}  & 4   & Max.\ 2-hop bridges sampled per undocumented pair \\
    \texttt{tier3\_candidates\_per\_pair}& 20  & Max.\ candidate labels offered for a Tier-3 pair \\
    \texttt{max\_serious/other\_effects} & 10 / 5 & Effects surfaced per analysis after ranking \\
    \bottomrule
  \end{tabular}
\end{table}

Because a documented drug pair can list dozens to hundreds of side effects
(a median of 53, up to 540, in Decagon), the effects of each pair are never
shown unranked. Every documented effect $e$ in the regimen receives a
transparent, additive score,
\begin{equation}
  \mathrm{score}(e) = w_r R(e) + w_i I(e) + w_m M(e) + w_s S(e)
                     + w_c\big[S(e)\wedge C(e)\big] + w_a A(e),
  \label{eq:score}
\end{equation}
where $R(e)$ is the number of regimen drug pairs that document $e$
(recurrence); $M(e)$, $S(e)$, and $C(e)$ are indicators for "absent from
either drug's individual side-effect list", "an EMA Important Medical
Event", and "additionally life-threatening"; $A(e)$ indicates that a pair
documenting $e$ also shares a protein target with a third regimen drug (the
Tier-2 trigger of Table~3.2); and $I(e)$ is an inverse pair-frequency term,
\begin{equation}
  I(e) = \log\!\left(\frac{N}{\max(n_e,\,1)}\right),
  \label{eq:idf}
\end{equation}
with $N$ the total number of distinct drug pairs documented in Decagon and
$n_e$ the number of pairs documenting $e$, so that effects specific to a
few pairs are weighted above ubiquitous ones. The six weights default to
$w_r{=}1.0$, $w_i{=}0.5$, $w_m{=}1.0$, $w_s{=}2.0$, $w_c{=}2.0$, $w_a{=}1.0$,
and are exposed through \texttt{Settings.ranking} so that they can be varied
in an ablation. Ranked effects determine both the order in which documented
effects are offered as evidence and, through $A(e)$, which pairs are
eligible for a Tier-2 prediction.

\subsection{Large Language Model Usage}
\label{sec:llm-usage}

Table~\ref{tab:llm} lists the request configuration. The LLM is accessed
through OpenRouter's\footnote{\url{https://openrouter.ai}} OpenAI-compatible
\texttt{/chat/completions} endpoint, which lets the model be swapped (via
the \texttt{LLM\_MODEL} environment variable) without changing application
code; the default model used in this project is Google's
\texttt{gemini-2.5-flash}\footnote{\url{https://ai.google.dev/gemini-api/docs/models}}.
The model is deliberately used as a constrained extraction-and-explanation
layer over evidence the pipeline has already retrieved, ranked, and bounded
(Section~\ref{sec:retrieval-impl}), rather than as an open-ended reasoner
over the raw graph: the request asks for \texttt{json\_object} output and
disables the model's own extended reasoning
(\texttt{reasoning: \{"effort": "none", "exclude": true\}}), which reduces
latency and cost without affecting prediction quality, since the biological
reasoning -- which proteins are shared, which bridges exist, which pairs
already have documented history -- is already encoded in the serialised
evidence rather than left for the model to derive on its own.

\begin{table}[t]
  \caption{LLM request configuration}
  \label{tab:llm}
  \centering
  \begin{tabular}{@{}ll@{}}
    \toprule
    Parameter & Value \\
    \midrule
    Endpoint         & OpenRouter, \texttt{/chat/completions} (OpenAI-compatible) \\
    Default model    & \texttt{google/gemini-2.5-flash} (overridable via \texttt{LLM\_MODEL}) \\
    \texttt{response\_format} & \texttt{json\_object} \\
    \texttt{reasoning}        & \texttt{\{"effort": "none", "exclude": true\}} \\
    \texttt{max\_tokens}      & 4000 \\
    Timeout          & 120\,s (15\,s connect) \\
    Retries          & up to 3, exponential back-off, on connection failure only \\
    \bottomrule
  \end{tabular}
\end{table}

The system prompt specifies the full output contract rather than leaving it
to free text: a JSON object with a \texttt{predictions} array whose items
carry \texttt{side\_effect}, \texttt{tier} ($1$--$3$), one or more
\texttt{evidence\_ids}, and a Vietnamese, at-most-three-sentence
\texttt{explanation} naming the drugs and, where relevant, the shared or
bridging protein targets. The prompt restates the tier semantics of
Table~3.2 as hard constraints on evidence use: a Tier-1 prediction must
reproduce a documented pairwise effect; a Tier-2 prediction must reproduce a
documented pairwise effect additionally supported by a target shared by at
least three regimen drugs; a Tier-3 prediction is only a hypothesis,
permitted solely when a retrieved bridge links the pair and that pair has
no documented effect, and its \texttt{side\_effect} label must be copied
verbatim from that bridge's \texttt{candidate\_side\_effects}. Output is
capped at 15 predictions, prioritising Tier~1, then Tier~2, then at most two
Tier-3 hypotheses, and the model is instructed to return an empty
\texttt{predictions} array rather than invent evidence when support is
insufficient.

\subsection{Grounding, Validation, and Fallback Behaviour}
\label{sec:grounding}

The raw response is parsed with a strict \texttt{pydantic} schema
(\texttt{ModelPrediction}/\texttt{PredictionList}); if that fails because
the model wrapped the JSON in a fenced code block or stopped mid-stream, a
tolerant recovery path scans the text with a streaming JSON decoder and
keeps every complete, schema-valid object found before the cut-off, rather
than discarding the whole response.

Parsing success does not make a prediction trustworthy: every item is
re-validated against the actually retrieved subgraph, independently of the
model's own claims. A prediction is kept only if every cited
\texttt{evidence\_id} exists, carries the declared tier, and (Tier~1/2)
names the same side-effect label on every cited item, or (Tier~3) lists
that label among the cited bridge's own candidates; the underlying graph
edges of each cited path are additionally checked against the loaded
target, PPI, and side-effect tables, and duplicate
(tier, label, evidence-set) predictions are removed. A prediction failing
any check is silently dropped rather than shown to the user. This two-stage
check -- schema validation, then graph grounding -- is what allows Tier-2
and Tier-3 output to be treated as traceable rather than as an unverified
model claim.

Because the LLM call can fail for several independent reasons (no API key
configured, a malformed or non-conforming JSON body, an HTTP error from the
provider, or a network failure), the pipeline treats all of them as a
single degraded mode: \texttt{llm\_status} is set to \texttt{"fallback"}
with a diagnostic message, and the request still returns the retrieved and
ranked subgraph, interaction graph, and coverage summary without raising an
exception, so a user still sees the underlying evidence even when no
generated explanation is available.

\end{document}
